\section{Lexical Structure}

This section defines the lexical tokens of OPQL. Whitespace and comments are recognized but discarded during lexical analysis.

\subsection{Keywords}

The following reserved keywords are recognized case-sensitively:

\begin{center}
\begin{tabular}{llll}
\toprule
Keyword & Alternatives & Keyword & Alternatives \\
\midrule
\term{PATTERN} & & \term{FILTER} & \\
\term{SUBJECTTO} & \term{ST} & \term{RETURN} & \\
\term{OCEL} & & \term{KEEP} & \\
\term{DISTINCT} & & \term{ORDERBY} & \\
\term{ASC} & \term{ASCENDING} & \term{DESC} & \term{DESCENDING} \\
\term{LIMIT} & & \term{BINNED} & \\
\term{WHEN} & & \term{AS} & \\
\term{MATERIALIZED} & & & \\
\term{True} & & \term{False} & \\
\term{None} & & \term{E} & \\
\term{O} & & \term{T} & \\
\term{D} & & \term{OR} & \\
\term{XOR} & & \term{AND} & \\
\term{NOT} & & & \\
\bottomrule
\end{tabular}
\end{center}

\subsection{Symbolic Names}

A symbolic name is an identifier used to bind values. It must start with a letter (uppercase or lowercase), followed by any sequence of letters, digits, or underscores.

\begin{framedgram}[Symbolic Names]
\begin{tabularx}{\linewidth}{l c X}
\nterm{symbolicName} & = & (\term{a}..\term{z} | \term{A}..\term{Z}) \{ \term{a}..\term{z} | \term{A}..\term{Z} | \term{0}..\term{9} | \term{\_} \} \\
\end{tabularx}
\end{framedgram}

The set of all valid symbolic names is denoted $\mathbb{U}_{sym}$.

\subsection{Strings}

A string literal is a sequence of characters enclosed in double quotation marks (\term{"}). The following escape sequences are recognized:

\begin{center}
\begin{tabular}{ll}
\toprule
Escape & Meaning \\
\midrule
\term{\textbackslash"} & Double quotation mark \\
\term{\textbackslash\textbackslash} & Backslash \\
\term{\textbackslash/} & Forward slash \\
\term{\textbackslash b} & Backspace \\
\term{\textbackslash f} & Form feed \\
\term{\textbackslash n} & Newline \\
\term{\textbackslash r} & Carriage return \\
\term{\textbackslash t} & Tab \\
\bottomrule
\end{tabular}
\end{center}

\subsection{Punctuation and Operators}

\begin{center}
\begin{tabular}{lll}
\toprule
Token & Name & Usage \\
\midrule
\term{,} & Comma & Separator \\
\term{:} & Colon & Type specifier \\
\term{(} \term{)} & Parentheses & Grouping, function calls \\
\term{[} \term{]} & Square brackets & Property lookup, relations \\
\term{*} & Asterisk & Wildcard, multiplication \\
\term{@} & At-sign & Timestamp qualifier for properties \\
\term{==} & Equal to & Comparison \\
\term{!=} & Not equal & Comparison \\
\term{<=} & Less than or equal & Comparison \\
\term{>=} & Greater than or equal & Comparison \\
\term{<} & Less than & Comparison, left-directed relation \\
\term{>} & Greater than & Comparison, right-directed relation \\
\term{+} & Plus & Addition, unary positive \\
\term{-} & Minus & Subtraction, unary negative, relation anchor \\
\term{/} & Slash & Division \\
\term{\%} & Percent & Modulo \\
\term{\textasciicircum} & Caret & Exponentiation \\
\term{->} & Right arrow & Right-directed relation \\
\term{<-} & Left arrow & Left-directed relation \\
\bottomrule
\end{tabular}
\end{center}

\subsection{Whitespace and Comments}

Whitespace (spaces, tabs, line breaks, and Unicode whitespace characters) is ignored except where needed to separate tokens.

Comments are supported in two forms:
\begin{itemize}
    \item Line comments: \term{//} followed by any characters until end of line
    \item Block comments: \term{/*} followed by any characters until \term{*/}
\end{itemize}
